\exnumber{04}
\extitle{UX I - Ex 5 - Énoncé du problème}
\exfiles{ex04/PDF, nomme Name\_FirstName\_DeliverableName\_Date\_VersionNumber}
\exauthorize{Ceux nommés dans l'exercice}
\makeheaderfiles

\textbf{Instructions :}

À partir des insights issus des analytics et des entretiens utilisateurs, commentez le périmètre des problem statements possibles.

\textbf{Livrables :}

\begin{itemize}
\item Sur un document PDF, rédigez un énoncé de problème qui résume le problème principal de Pierre lorsqu'il court.
\item Déposez le PDF sur Github, nommé ainsi : « Name\_FirstName\_DeliverableName\_Date\_VersionNumber ». Exemple : « Doe\_John\_MidFiPrototype\_05242024\_V1 ».
\item Ajoutez un titre dans le document.
\item Soignez l'aspect global du document. Il doit être clair, simple et lisible. Vous pouvez regarder canva.com, sans surcharger le dessin de détails.
\end{itemize}

\textbf{Pistes :}

\begin{itemize}
\item À ce stade, il est important de se concentrer sur le problème, pas sur la solution ! Vous pourrez vous attaquer au problème après l'avoir (re)défini.
\item Plusieurs énoncés de problème peuvent être possibles pour chaque persona et parcours utilisateur, mais pour cet exercice, nous vous orientons vers un seul.
\item Exemples d'énoncés de problème corrects :
\begin{itemize}
\item Comment pourrions-nous aider Pierre (ou votre persona) à mesurer et préparer son hydratation avant une course ?
\item Comment pourrions-nous nous assurer que Pierre ne se déshydrate pas ?
\item Comment pourrions-nous nous assurer que Pierre boit les gobelets proposés pendant la course ?
\item Comment pourrions-nous fournir à Pierre une boisson qui contient tout ce dont il a besoin pour éviter la déshydratation ?
\item Comment pourrions-nous nous assurer que Pierre ne se soûle pas la veille de la course ?
\end{itemize}
\item Exemples d'énoncés de problème incorrects :
\begin{itemize}
\item Comment pourrions-nous nous assurer que Pierre ne se déshydrate pas ?
\begin{itemize}
\item Trop large  --- > Le problème est hors du périmètre de la course à pied, donc les solutions ne peuvent pas être intégrées dans le contexte spécifique de la course à pied.
\end{itemize}
\item Comment pourrions-nous faire en sorte que Pierre boive les gobelets proposés pendant la course ?
\begin{itemize}
\item Trop étroit  --- > Il existe plein d'options pour éviter la déshydratation sans boire les gobelets proposés pendant la course.
\end{itemize}
\item Comment pourrions-nous fournir à Pierre une boisson qui contient tout ce dont il a besoin pour éviter la déshydratation ?
\begin{itemize}
\item Centré sur la solution, pas sur le problème  --- > L'énoncé du problème est censé reformuler le brief et prendre en compte les insights de la phase précédente.
\end{itemize}
\item Comment faire pour que Pierre ne se soûle pas la veille de la course ?
\begin{itemize}
\item Hors périmètre  --- > Il ne s'agit ni de déshydratation, ni de course à pied.
\end{itemize}
\end{itemize}
\end{itemize}

\textbf{Ressources :}

\begin{itemize}
\item Bloqué sur un problème ? Demandez simplement How Might We \url{https://relab.academy/design-thinking/stuck-on-a-problem-just-ask-how-might-we/}
\item Problem statement \url{https://www.indeed.com/career-advice/career-development/what-is-a-problem-statement}
\item Utiliser les questions « How Might We » pour idéer sur les bons problèmes \url{https://www.nngroup.com/articles/how-might-we-questions/}
\end{itemize}

\textbf{Citation :}

« Si j'avais une heure pour résoudre un problème, je passerais 55 minutes à réfléchir au problème et 5 minutes à réfléchir aux solutions. » Albert Einstein

